\section{Problem formulation and notation}
\label{sec:setup}

\subsection{Observations, coordinates, and admissible partitions}
Let $\{(x_i,y_i,d_i)\}_{i=1}^n$ be observations with ordered design points
$x_1<\cdots<x_n$. The response $y_i$ may be continuous or discrete; $d_i$
contains known observation-model quantities, such as a Poisson exposure, a
Binomial trial count, or a Gaussian precision. A separate indicator $o_i\in
\{0,1\}$ specifies whether the response is observed. Missing outcomes do not
contribute likelihood factors. We write $y_{\mathrm{obs}}$ for the observed
responses. The default analysis uses ordinary sampling likelihoods. When
likelihood powers $a_i\geq0$ are deliberately introduced, they define a
\emph{generalized Bayesian} calculation and are not interpreted as exposure,
trial count, or physical spacing.

Candidate partition edges are $u_0<u_1<\cdots<u_n$, with interior edge $u_i$
separating observations $i$ and $i+1$. For irregular observations, an admissible
choice is $u_i=(x_i+x_{i+1})/2$ for $1\leq i<n$, together with prespecified
endpoints $u_0<x_1$ and $u_n>x_n$. Where the observation domain has known
physical limits, use those limits instead. Choosing these edges is part of the
model specification. On an index grid, use $u_i=i$. These physical edges are
distinct from the observed positions $x_i$ and from the family descriptor $d_i$.

A partition into $k$ nonempty contiguous segments has ordered boundary vector
\begin{equation}
 t=(t_0,\ldots,t_k),\qquad 0=t_0<t_1<\cdots<t_k=n,
 \qquad I_q(t)=\{t_{q-1}+1,\ldots,t_q\}.
 \label{eq:setup-partition}
\end{equation}
The candidate block $(i,j]$ denotes observations $i+1,\ldots,j$ and has
physical length $u_j-u_i$. Let $\mathcal T_k$ be the set of partitions satisfying
the declared candidate-boundary and minimum-length restrictions. These
restrictions are set independently of the responses. Write $1\leq k\leq
k_{\max}$ for the admissible segment counts.

\subsection{Hierarchical sampling model and prior}
Within segment $q$, observations are conditionally independent given the
segment parameter $\theta_q$:
\begin{equation}
 p(y_{\mathrm{obs}}\mid t,k,\theta_{1:k},d)
 =\prod_{q=1}^k\ \prod_{i\in I_q(t):o_i=1}
 p(y_i\mid\theta_q,d_i),
 \qquad \theta_q\overset{\mathrm{ind}}\sim\pi_0.
 \label{eq:setup-likelihood}
\end{equation}
The segment prior $\pi_0$ is proper and does not depend on the partition unless
explicitly specified. The prior on the count and the conditional prior on the
partition are $p(k)$ and $p(t\mid k,u)$, respectively. For the block-local
factorization used below, write
\begin{equation}
 p(t\mid k,u)=\frac{1}{C_k(u)}\prod_{q=1}^k
 \gamma_u(t_{q-1},t_q),\qquad
 C_k(u)=\sum_{t\in\mathcal T_k}\prod_{q=1}^k
 \gamma_u(t_{q-1},t_q).
 \label{eq:setup-prior}
\end{equation}
The definition assumes $0<C_k(u)<\infty$. Segment cohesion and
candidate-boundary hazards are two components of $\gamma_u$ and are defined in
Section~\ref{sec:model}. For the uniform conditional prior, $\gamma_u\equiv1$
on $\mathcal T_k$ and $C_k=|\mathcal T_k|$; in the unrestricted case,
$C_k=\binom{n-1}{k-1}$.

\subsection{Inference and decision targets}
Define the integrated likelihood of a candidate block by
\begin{equation}
 A^{(0)}_{ij}=\int\!\pi_0(\theta)
 \prod_{r=i+1}^j p(y_r\mid\theta,d_r)^{o_r}\,\mathrm d\theta.
 \label{eq:setup-block}
\end{equation}
Here a missing response contributes the multiplicative identity one; in
particular, a block containing no observed responses has $A^{(0)}_{ij}=1$ under
a proper prior. The count-conditional marginal likelihood and posterior are
\begin{equation}
 p(y_{\mathrm{obs}}\mid k,u)
 =\frac{1}{C_k(u)}\sum_{t\in\mathcal T_k}
 \prod_{q=1}^k \gamma_u(t_{q-1},t_q)
 A^{(0)}_{t_{q-1},t_q},\qquad
 p(k\mid y_{\mathrm{obs}},u)\propto p(k)\,p(y_{\mathrm{obs}}\mid k,u).
 \label{eq:setup-evidence}
\end{equation}
The inferential targets are the full count posterior, ordered-boundary
marginals $p(t_q=h\mid y,k)$, boundary-event probabilities
$p(b_h=1\mid y,k)$, posterior moments of a declared segment signal
$m_\star(\theta)$, and the posterior predictive distribution. The
\emph{marginal-MAP count} maximizes $p(k\mid y)$; a \emph{conditional joint MAP
partition} maximizes $p(t\mid y,\widehat k)$; the \emph{global joint MAP pair}
maximizes $p(k,t\mid y)$. These are different posterior summaries and can
select different partitions. Exact formulas and algorithms are given in
Section~\ref{sec:exact-inference}.

\subsection{Extension to related sequences}
For $S$ aligned sequences, let $y^{(s)}_i$ and $\theta^{(s)}_q$ denote the
subject response and its segment parameter. In the shared-boundary model,
$t$ is common to all $S$ sequences while the $\theta^{(s)}_q$ are independent
given $t$. Known groups have distinct group-specific partitions. For unknown
group labels, Section~\ref{sec:extensions} defines a separate finite-template
score criterion and its minorization--maximization procedure. The latter is
not a posterior distribution over a normalized finite mixture unless an
additional probabilistic model is supplied.

\begin{table}[!htbp]
\centering\small
\caption{Notation used throughout the manuscript.}
\label{tab:notation-paper}
\begin{tabularx}{0.96\textwidth}{@{}L{0.22\textwidth}Y L{0.19\textwidth}Y@{}}
\toprule
Symbol & Definition & Symbol & Definition \\
\midrule
$x_i$, $u_i$ & Observed position; candidate partition edge & $y_i,d_i,o_i$ & Response, family descriptor, observation mask \\
$n,S,G$ & Observations, sequences, latent groups & $k,k_{\max}$ & Segment count and imposed upper bound \\
$t,\mathcal T_k$ & Partition and its admissible set & $I_q(t)$ & Index set of segment $q$ \\
$\theta_q,\pi_0$ & Segment parameter and proper prior & $m_\star(\theta)$ & Reported segment-level signal functional \\
$A^{(0)}_{ij}$ & Integrated block likelihood & $\gamma_u(i,j)$ & Local partition-prior factor \\
$C_k(u)$ & Conditional partition-prior normalizer & $L_{k,j},R_{k,i}$ & Forward and backward partition sums \\
$M_{k,j}$ & Max-sum optimal prefix score & $p(b_h=1\mid y,k)$ & Conditional boundary-event probability \\
\bottomrule
\end{tabularx}
\end{table}
